\documentclass{article}
\usepackage{iclr2026_conference,times}
\usepackage{amsmath,amssymb}
\usepackage{graphicx}
\usepackage{booktabs}
\usepackage{hyperref}
\usepackage{url}

\title{Grading the Diagnostic Conversation:\\
Causal-Graph-Grounded, Execution-Free Evaluation\\
of Debugging Agents on Real Support Threads}

\author{Anonymous}

\newcommand{\fix}[1]{\textsf{#1}}

\begin{document}
\maketitle

\begin{abstract}
An agent that debugs \emph{with a person} must elicit the evidence a fix
depends on, avoid the attempts already known to fail, and ground its
conclusion in what the user actually reported. Existing multi-turn
benchmarks measure none of this directly: they either require a
reproducible execution environment---which excludes precisely the
environment-bound problems users bring---or condition a simulated user on
the resolved transcript, leaking the answer into the dialogue. We
annotate 229 resolved public support threads as causal graphs whose nodes
are (system-state, information-state) pairs and whose edges are
clarifications, solutions, and attempts the thread itself falsified. The
same graph constrains a leak-safe user simulator, grounds an
execution-free judge, and yields metrics that need no model at all. Four
frontier models separate into two tiers 0.32--0.35 apart, and the split
reproduces under a judge from a different model family and under
judge-free counts alone. Two results are less comfortable: on staying
inside the evidence the strongest models score 0.64--0.71 and the weaker
pair 0.05--0.09, and when a single authored answer is swapped for a
plausible alternative, the strongest model stops being able to close the
case at all 72\% of the time---it does not adapt, it derails. We also
report what auditing the harness cost: before repair, only 7--11\% of the
benchmark's own intervention events reflected genuine agent failure.
\end{abstract}

\section{Introduction}

Debugging with a user is an elicitation problem before it is a repair
problem. The information a fix depends on is usually not in the report:
in our corpus 79\% of the facts a resolution turns on had to be asked
for or measured. An agent that proposes early proposes blind.

Benchmarks for multi-turn software agents take one of two routes, and
both lose the thing we want to measure. Execution-grounded suites need a
containerised reproduction, which systematically excludes the
environment-bound faults real users report---drivers, devices, OS
integration, account state. Simulated-user suites avoid that cost by
scripting the user, and typically condition the simulator on the
resolved thread, so the user knows the answer and the dialogue is no
longer a diagnosis.

We take a third route. Each case is a resolved public thread annotated as
a causal graph: nodes are (system-state, information-state) pairs, edges
are clarifications the user can answer, solutions with required elements,
and attempts this thread tried and falsified. The graph is an answer key,
not a transcript to replay. The same object then does three jobs: it
bounds what the simulated user may say (it speaks only from its current
node), it gives an execution-free judge something to score against, and
it makes several metrics deterministic---derivable from the transcript
with no model in the loop.

Our contributions:

\begin{enumerate}
\item A DSL and a corpus: 229 released cases over 75 projects, each
      machine-validated and passed by a two-stage model review, with the
      falsified attempts annotated as first-class structure.
\item An evaluation protocol whose central comparison does not depend on
      an LLM judge, and a measurement of what the judge adds where it is
      used.
\item Results on four frontier models, including a counterfactual
      intervention study that the corpus design uniquely enables.
\item An audit of the harness itself. We report the defects we found in
      our own instrument, and what each was worth, because several of
      them would have been read as model behaviour.
\end{enumerate}

\section{Related Work}
\label{sec:related}

\textbf{Execution-grounded agent benchmarks.} SWE-bench and its
descendants score a patch against a test suite. This is decisive where
it applies and inapplicable where reproduction is the hard part;
CAB's own limitations section documents the resulting selection bias.

\textbf{Simulated-user benchmarks.} The $\tau$-family and CAB place an
LLM in the user's seat. Fidelity is the open question: a simulator
conditioned on the resolution can leak it, and the leak is invisible in
aggregate scores. We measure this directly in Section~\ref{sec:leak}.

\textbf{Clarification and elicitation.} A line of work scores whether a
model asks a question at all. We score whether it obtained the specific
evidence its own proposal required, which the graph makes checkable.

\section{Task Formulation}
\label{sec:task}

\begin{figure}[t]
\centering
\includegraphics[width=\textwidth]{f1_worked_case.pdf}
\caption{One case as a causal graph. Solid: the diagnostic path. Dotted:
the attempt this thread actually made and falsified, together with the
aftermath state it leads to. An agent may take the falsified branch---the
case's satisfaction conditions forbid \emph{reporting} it as the fix.}
\label{fig:case}
\end{figure}

Figure~\ref{fig:case} shows a complete case. A case is a directed
graph over states. A state pairs a system state with
an information state: what is true of the system, and what has been
established in the conversation. Edges are of three kinds.
\emph{Clarification} edges carry one or more facts the user can supply,
each labelled by how hard it is to obtain: \textbf{L1} any user states it
unprompted, \textbf{L2} inferable from evidence already surfaced,
\textbf{L3} must be asked for or measured. \emph{Solution} edges carry an
intent, a concrete example, and the elements a reply must contain to
count as that fix, plus the information the fix depends on.
\emph{Known-blind} edges are solutions the thread tried and falsified;
they lead to aftermath states from which the investigation continues.

Two derived layers are generated rather than authored: rollback edges out
of every aftermath state, and shortcut copies of a solution edge planted
at earlier states that share its system state, which let an agent skip
evidence---and let the judge grade the skip.

\textbf{The answer-key principle.} The out-edges of a state are what
\emph{could} happen there, not a replay of what the historical
participants did. This is what makes a counterfactual intervention
meaningful (Section~\ref{sec:cf}).

\section{Corpus}
\label{sec:corpus}

\begin{figure}[t]
\centering
\includegraphics[width=\textwidth]{f8_corpus.pdf}
\caption{Corpus shape. Most evidence a resolution depends on has to be
asked for: 79\% of clarifications are L3.}
\label{fig:corpus}
\end{figure}

The corpus (Figure~\ref{fig:corpus}) holds 229 released cases across 75
projects, no project exceeding 4\% of the corpus, drawn from Mozilla Bugzilla, PostgreSQL, and permissive-license
GitHub repositories, threads created 2021--2026 (median 37 messages).
Cases are drafted by a model from the thread and gated by schema
validators and structural lints; a two-stage review then files findings
and independently verifies them. Single-pass review refuses roughly 95\%
of cases, mostly on claims its own evidence does not support;
verification is what makes model review usable as a release gate. Of 29
cases blocked in the first round, 25 passed after redrafting under
hardened rules---their defects came from the generation prompt, not from
the threads.

\textbf{Not recall-solvable.} Handed the opening report alone, with no
opportunity to ask anything, a model names both mechanism and fix class
in 2.2\% of cases (partial credit in a further 3.5\%), with no gradient
by issue year. The benchmark measures interactive diagnosis, not
memorisation.

\section{Evaluation Protocol}
\label{sec:protocol}

\textbf{The simulated user} speaks only from its current graph position:
that state's visible symptoms, authored answers to clarifications the
agent actually asked, and---when a stall-insurance mechanism fires---a
reveal that is recorded as such. It is never shown the resolution.

\textbf{Matching.} Each agent turn is read for both acts. A question is
matched against the case's clarifications; a proposal against the current
state's solution edges, producing a \emph{solution call} that records
which required information the agent held at that moment, by level.

\textbf{Metrics.} Case outcomes (\fix{terminal\_resolved}, forced walk,
ran out), turns, forced reveals per case, and premature proposals per
case are deterministic given the transcript. On top of these, four LLM
rubrics---asks for evidence before proposing, stays inside the evidence,
accounts for the fault, recovers from a failed step---combine with the
graph-supplied grounding rate into a single \fix{grade}.

\textbf{Ruling on a difference.} Per-case grades are paired across arms.
A difference is called real only when $|t| > 2.6$ \emph{and} a two-sided
sign test gives $p < 0.01$. This replaced an earlier rule of ``at least
twice the run-to-run drift'', which ignores sample size and consequently
misruled in both directions.

\section{Experiments}
\label{sec:exp}

\begin{table}[t]
\caption{Four models over all 229 cases, 30 turns.}
\label{tab:main}
\centering
\begin{tabular}{lrrrrrr}
\toprule
model & grade & resolved & forced walk & ran out & turns & reveals/case \\
\midrule
gpt-5.6  & \textbf{0.624} & 66 (29\%) & 110 (48\%) & 45 (20\%) & 21 & 3.1 \\
gpt-5.5  & \textbf{0.574} & 71 (31\%) & 103 (45\%) & 47 (21\%) & 20 & 2.8 \\
GLM-5.1  & \textbf{0.303} & 59 (26\%) & 115 (50\%) & 49 (21\%) & 22 & 3.7 \\
Kimi-2.5 & \textbf{0.279} & 52 (23\%) & 103 (45\%) & 70 (31\%) & 21 & 4.0 \\
\bottomrule
\end{tabular}
\end{table}

Table~\ref{tab:main} gives the main result. Two tiers separate by
0.32--0.35 ($t = 32$); within the stronger tier gpt-5.6 leads gpt-5.5 by
0.050 ($t = 5.3$), and between GLM-5.1 and Kimi-2.5 there is nothing to
call. Fewer than a third of cases end in a fix the agent earned, and
roughly half reach a terminal state only because the insurance walked
them there.

\begin{figure}[t]
\centering
\includegraphics[width=0.86\textwidth]{f2_rubrics.pdf}
\caption{The tiers separate on all four rubrics and never cross. The
deepest gap is staying inside the evidence.}
\label{fig:rubrics}
\end{figure}

\subsection{Where the weaker tier loses}

Figure~\ref{fig:rubrics} localises the gap. It is not confined to any one
rubric, and the largest is \emph{staying inside the evidence}: the GPT
pair scores 0.64--0.71, the other two 0.05--0.09---asserting what their
conversations do not support in nearly every case.

\subsection{The central comparison does not need the judge}

Two per-case counts involve no model scoring: forced reveals, and
proposals made from a state whose graph offers no solution edge. Paired
across the same 229 cases, every cross-tier comparison separates
(premature proposals $+1.41$ and $+1.64$, $t = 6.7$ and $6.9$) and every
within-tier comparison does not---the same partition
Table~\ref{tab:main} reports, obtained by counting rather than scoring.

\subsection{Counterfactual interventions}
\label{sec:cf}

The corpus authors, for clarifications that mattered, an alternative
answer the reporter could plausibly have given, marked with whether the
right fix changes as a result. Replaying a case with one answer swapped
asks whether the agent is conditioning on what the user said. Drawing 60
interventions only from cases whose baseline run closes unaided gives
Table~\ref{tab:cf}:

\begin{table}[h]
\caption{One authored answer swapped, gpt-5.6, $n=60$.}
\label{tab:cf}
\centering
\begin{tabular}{lrr}
\toprule
 & fix should change & fix should not change \\
\midrule
proposed something different & \textbf{5 (17\%)} & 2 (7\%) \\
proposed the same thing & 5 (17\%) & \textbf{5 (17\%)} \\
never reached a proposal of its own & 20 (67\%) & 23 (77\%) \\
\bottomrule
\end{tabular}
\end{table}

Changing one answer stops the agent closing the case at all 72\% of the
time. It does not switch to a different fix; it stops converging---and it
derails \emph{more} often when the change was marked as not affecting the
diagnosis. On the 17 cases that stayed on the rails, five of ten
interventions that should have changed the fix did.

\subsection{Does the result survive a different judge?}

The primary judge is from the same family as one of the models it ranks.
Re-judging the same transcripts with GLM-5.1 under identical rubric
definitions gives gpt-5.6 over Kimi-2.5 at $+0.189$ ($t = 17.7$) against
the primary judge's $+0.345$ ($t = 31.7$). Absolute levels are
judge-specific; direction, significance and ordering hold.

\subsection{Leakage}
\label{sec:leak}

Handing the simulator the resolved thread---the transcript-conditioned
design---produces leaked turns that a lexical screen detects (3 in 916
against 0 in 913 under our production profile). The grade inflation this
causes is \emph{not} established at our sample size ($+0.037$, $t = 1.8$,
$p = 0.19$, $n = 48$), and we report it as such.

\section{Auditing the Instrument}
\label{sec:audit}

\begin{figure}[t]
\centering
\includegraphics[width=0.86\textwidth]{f6_all_comparisons.pdf}
\caption{Every comparison, with a 99\% interval. The final row is the
same configuration run twice: interventions that changed nothing sit
where re-running sits.}
\label{fig:all}
\end{figure}

Figure~\ref{fig:all} places every comparison we ran on one axis, with
the same configuration run twice as the bottom row. Building this
benchmark produced findings about the benchmark. We report
them because several would otherwise have been read as model behaviour.

\textbf{Most intervention events were our fault.} Before repair, 44\% of
cases in every row triggered a forced reveal at a rate that barely varied
across models---a signal belonging to the harness. Decomposed, only
7--11\% reflected an agent failing a modelled path; the rest were a
matcher that could not represent a turn containing both a question and a
proposal, a simulator that fabricated negative results for fixes the case
never modelled, and a canonical-path search that refused the falsified
edges most threads pass through, killing half of every run at a median of
turn 5.

\textbf{Two published findings were rubric artefacts.} A rubric asking
whether the agent ``proactively gathered info'' scored every model
0.92--0.97 and supported a conclusion that elicitation was saturated.
A judge-independent count contradicted it the whole time---premature
proposals per case run 2.36 to 4.00 across the four models---and under a
rubric that asks whether the evidence was in hand \emph{before} the
proposal, the scores track that count at Spearman $-0.80$. Separately, an
under-specified rubric made two judges score different constructs and
produced an apparent finding that the tier gap belonged to the judge; it
did not survive the rubrics being written down.

\textbf{What we now check on every run.} Ten offline simulator
invariants; a row-integrity reconciliation of transcripts against metrics
against judgments, flagging when the missing cases are disproportionately
successes; a validity gate on empty replies; a leak screen; and a rubric
sanity check against the behavioural counts.

\section{Limitations}

The corpus is machine-drafted and machine-reviewed, with the audit trail
published. \textbf{The judge was never validated against a human.} Two
checks stand in its place---a judge from another family reproduces the
ranking, and the rubrics track a judge-independent count---and neither is
equivalent; two models agreeing is not correctness. Mitigating this, the
tier result is also established on judge-free counts alone, so the
central comparison does not rest on the judge, though the rubric-level
readings do. Per-case grades are unstable across identical runs (mean
absolute drift 0.10--0.12), so no claim rests on an individual case.
\fix{reached\_terminal} saturates once cases use their full budget and is
reported as harness health rather than as an agent metric. Threads are
English-only and from open-source projects.

\section{Release}

Code under Apache-2.0; the curation layer under CC BY 4.0; thread content
remains its authors', collected from permissive sources with usernames
pseudonymised. The 55 drafts that review blocked ship alongside the 229
released cases with their findings, and the released set is exactly those
carrying \fix{hitl\_reviewed}.

\end{document}
